\lysh{38916}{4}

\begin{alist}
\item Από τη γραφική παράσταση που μοντελοποιεί το άλμα παρατηρούμε:

\begin{center}
\begin{tabular}{|c|c|c|c|c|c|}

χρόνος (sec) & $60$ & $65$ & $74$ & $78$ & $84$ \\ 
ύψος (m) & $38$ & $42$ & $42$ & $40$ & $40$ \\ 
\end{tabular}
\end{center}

\item Όπως φαίνεται συνδυαστικά από τον πίνακα και τη γραφική παράσταση:
\begin{itemize}
\item για τη χρονική στιγμή $0\text{ sec}$, ο άνθρωπος πηδάει από ύψος $80\text{ m}$,
\item το μικρότερο ύψος στο οποίο φτάνει είναι τα $13\text{ m}$,
\item η χρονική στιγμή κατά την οποία φτάνει σε ύψος $13\text{ m}$ είναι τα $11\text{ sec}$.
\end{itemize}

\item Τα πρώτα $40\text{ sec}$ ο άνθρωπος φτάνει σε ύψος $40\text{ m}$ σε τέσσερις χρονικές στιγμές:
\begin{itemize}
\item στα $6\text{ sec}$ και στα $40\text{ sec}$ (όπως προκύπτει από τον πίνακα),
\item στα $17\text{ sec}$ περίπου και στα $29\text{ sec}$ περίπου (όπως προκύπτει από τη γραφική παράσταση).
\end{itemize}

\item Μετά τα πρώτα $40\text{ sec}$, η μικρότερη τιμή του ύψους στο οποίο φτάνει ο άνθρωπος είναι περίπου $38\text{ m}$ και η μεγαλύτερη τιμή είναι $46\text{ m}$. Επομένως, η διαφορά ύψους χαμηλότερου και ψηλότερου σημείου της κίνησης του ανθρώπου είναι περίπου $8\text{ m}$, δηλαδή ο ισχυρισμός είναι λανθασμένος.

\item Όπως προκύπτει από τη γραφική παράσταση, το ύψος στο οποίο βρίσκεται ο άνθρωπος όσο περνά ο χρόνος, για παράδειγμα τις χρονικές στιγμές μετά τα $100\text{ sec}$, κυμαίνεται μεταξύ $40\text{ m}$ και $42\text{ m}$. Επομένως, η διαφορά ύψους χαμηλότερου και ψηλότερου σημείου της κίνησης του ανθρώπου, όσο περνά ο χρόνος, μειώνεται και μένει σταθερά κάτω από τα $2\text{ m}$.
\end{alist}

